% Conclusions: restate the contribution and the forward path. Instruction only.

[PROCESSING INSTRUCTION (Conclusions). Write one or two tight paragraphs,
no table. Restate the headline finding in the author's voice: an
autonomous Claude Code Opus 4.7 (1M context) Max run built, executed, and
documented a standard-library oncology automation platform under strict
assurance, and the evidence supports the thesis directly, because the
pipeline produced a complete deliverable in under a millisecond while the
VVUQ gate accepted only 1 of 6 candidates and blocked 5. Conclude that the
assurance process, not the generation, is the substantial and decision
bearing part, and that holding VVUQ higher than code generation is what
makes physical AI oncology trial deliverables faster, less expensive, and
more rigorous than conventional verification.

End on the forward path: the same gate scales to the Stage 2 2030 PDAC
analog and, with an independent validation reference and the real input
corpus wired in, to oncology clinical trials broadly, so the field can
check these assurance steps off once rather than repeatedly, with patient
benefit preserved throughout.

SOURCE FILES: execution/README.md (Thesis and the execution results
summary) and the Results and Discussion of this paper. Do not introduce
new numbers that were not established earlier.

HOW TO PROCESS: this is a synthesis, not a summary list; connect back to
the Introduction's framing of Verification, Validation, and Uncertainty
Quantification. Keep one citation at most, for example
\cite{repo-cancer-automated}; avoid a citation pile-up in the conclusion.]
